% Lösungen Potenz- und Exponentialfunktionen, Wachstum und Zerfall (zusammenbau v0.4, Heft)
\einheitenkopf{Potenz- und Exponentialfunktionen, Wachstum und Zerfall}

\erg{31}{a) der Betrag \quad b) $30$; $30$; $30$ – gleich \quad c) exponentiell – jeder Wert ist das $0{,}8$-Fache des vorigen}
\erg{32}{a) Exponentiell: jeder Wert ist das $1{,}06$-Fache des vorigen; die Zuwächse ($1\,500$, dann $1\,590$) werden größer. \quad b) Exponentiell: jeder Umsatz ist das $1{,}12$-Fache des vorigen, die Zuwächse in Euro wachsen ($18\,000$, dann $20\,160$). \quad c) Die Menge nimmt jede Stunde um denselben Anteil ab. – Faktor $0{,}85$ je Stunde; die Abnahmen in mg ($6$; $5{,}1$; $4{,}33$) werden kleiner. \quad d) Der Druck sinkt jeden Tag um denselben Prozentsatz. – jeder Wert ist das $0{,}9$-Fache des vorigen; die Abnahmen ($0{,}25$; $0{,}225$; $0{,}2025$) werden kleiner.}
\erg{33}{a) $h$ – nur er geht durch $(0 | 0)$ \quad b) $f$ Kurve; $g$ Gerade; $h$ Kurve \quad c) $f$ – fällt und wird flacher, weil die Abnahme in Euro kleiner wird}
\erg{34}{a) $f$ – Start bei $80$ und gekrümmt. $g$ nicht: Gerade, also jedes Jahr derselbe Betrag. $h$ nicht: beginnt bei $0$, der Tarif aber bei $80$. \quad b) $u$ – Start bei $40$, gekrümmt. $v$ nicht: Gerade, fester Zuwachs. $w$ nicht: Start im Ursprung, die Stadt hat aber schon $40$ Tausend Einwohner. \quad c) $f$ passt: Start bei $45$, Kurve. $g$ passt nicht: Gerade, also jede Woche gleich viel. $h$ passt nicht: Start bei $0$, das Kalb wiegt aber $45\,\mathrm{kg}$. \quad d) $u$ passt: Start bei $600$, gekrümmt. $v$ passt nicht: Gerade, gleicher Betrag je Jahr. $w$ passt nicht: beginnt bei $0$.}
\erg{35}{a) $350$ \quad b) Punkte $(0 | 100)$, $(1 | 120)$, $(2 | 144)$, $(3 | 173)$, $(4 | 207)$ \quad c) Punkte $(0 | 125)$, $(1 | 150)$, $(2 | 180)$, $(3 | 216)$ – $125$ liegt genau in der Mitte zwischen $100$ und $150$}
\erg{36}{a) Punkte $(0 | 56)$, $(1 | 84)$, $(2 | 126)$, $(3 | 189)$, $(4 | 284)$ \quad b) Punkte $(0 | 500)$, $(1 | 550)$, $(2 | 605)$, $(3 | 666)$, $(4 | 732)$ \quad c) waagerecht $1\,\mathrm{km}$ je Kästchen, senkrecht $100\,\mathrm{hPa}$ je Kästchen; Punkte eingetragen und fallende Kurve \quad d) waagerecht $1\,\mathrm{m}$ je Kästchen, senkrecht $100\,\mathrm{lx}$ je Kästchen; Punkte eingetragen, fallende Kurve \quad e) waagerecht $1\,\mathrm{h}$ je Kästchen, senkrecht $1\,\mathrm{mg}$ je Kästchen; Punkte eingetragen und fallende Kurve \quad f) waagerecht $1$ Tag je Kästchen, senkrecht $10\,\mathrm{g}$ je Kästchen; Punkte eingetragen, fallende Kurve}
\erg{37}{a) mehr – der Faktor ist größer als eins \quad b) $1{,}06$ \quad c) $0{,}94$}
\erg{38}{a) $80 : 64 = 1{,}25$; $100 : 80 = 1{,}25$; $125 : 100 = 1{,}25$ – alle Quotienten sind $1{,}25$. \quad b) $240 : 120 = 2$; $480 : 240 = 2$; $960 : 480 = 2$ – der Faktor ist überall $2$. \quad c) $250\,000 \cdot 1{,}05 = 262\,500\,€$ (oder $12\,500 : 250\,000 = 0{,}05 = 5\,\%$) \quad d) $8\,000 \cdot 1{,}03 = 8\,240$ (oder $240 : 8\,000 = 0{,}03 = 3\,\%$)}
\erg{39}{a) ausgelassen \quad b) $250 \cdot 1{,}15 = 287{,}5$ \quad c) $300$ – der Startwert steht im Jahr $0$, er wird nicht noch einmal mit $1{,}1$ multipliziert}
\erg{40}{a) $80{,}00$ und $86{,}15$ ($2026$ ist der Anfangswert; $84{,}05 \cdot 1{,}025 \approx 86{,}15$) \quad b) $240{,}00$ und $262{,}26$ (Anfangswert $2025$; $254{,}62 \cdot 1{,}03 \approx 262{,}26$) \quad c) $36{,}0$ und $40{,}5$ ($2015$: Startwert; $38{,}9 \cdot 1{,}04 \approx 40{,}5$) \quad d) $120{,}0$ und $127{,}3$ (Startwert; $124{,}8 \cdot 1{,}02 \approx 127{,}3$) \quad e) $6{,}40\,\mathrm{mg}$ und $4\,\mathrm{h}$: der Wert bei $1\,\mathrm{h}$ und die Zeitangabe, denn $8 \cdot 0{,}8^4 = 3{,}28$ ($5{,}12 \cdot 0{,}8 = 4{,}10 \ne 3{,}28$) \quad f) $60\,\mathrm{mg}$ und $3$: der Wert bei Tag $1$ und die Zeitangabe Tag $3$, denn $80 \cdot 0{,}75^3 = 33{,}75$ \quad g) $4{,}000$; $5{,}018$; $7{,}049$ (Woche $0$ Startwert; $4{,}48 \cdot 1{,}12$; $4 \cdot 1{,}12^5$) \quad h) $6{,}000$; $6{,}998$; $9{,}521$ (Startwert; $6{,}48 \cdot 1{,}08$; $6 \cdot 1{,}08^6$) \quad i) $620$ und $371$ ($704 \cdot 0{,}88 \approx 620$; $545 \cdot 0{,}88^3 \approx 371$) \quad j) $723$ und $321$ ($850 \cdot 0{,}85 \approx 723$; $522 \cdot 0{,}85^3 \approx 321$)}
\erg{41}{a) $t = 5$ – $2016$ ist der Schritt null \quad b) $3\,000$ und $1{,}04$ (Anfangswert und Faktor) \quad c) $E(t) = 3\,000 \cdot 1{,}04^t$}
\erg{42}{a) $S(t) = 80 \cdot 1{,}025^t$; $2038$: $t = 12$, $S(12) = 80 \cdot 1{,}025^{12} \approx 107{,}59\,€$ \quad b) $G(t) = 240 \cdot 1{,}03^t$; $2035$: $t = 10$, $G(10) = 240 \cdot 1{,}03^{10} \approx 322{,}54\,€$ \quad c) $y = 800 \cdot 0{,}88^x$ – Anfangswert mal Faktor hoch $x$, Faktor $0{,}88$ für die Abnahme \quad d) $y = 1\,000 \cdot 0{,}85^x$ – Faktor $0{,}85$ je Meter, $x$ im Exponenten \quad e) $4$: Masse bei der Geburt in kg; $1{,}12$: Wachstumsfaktor, $100\,\% + 12\,\%$ je Woche; $x$: Anzahl der Wochen seit der Geburt \quad f) $350$: Besucher zu Beginn; $1{,}25$: Wachstumsfaktor, jeden Monat $25\,\%$ mehr; $x$: Anzahl der Monate seit Beginn}
\erg{43}{a) $1{,}04^5 \approx 1{,}2167$ \quad b) $440$; $484$; $532{,}4$ \quad c) $t = 6$: $2\,500 \cdot 1{,}06^6 \approx 3\,546$}
\erg{44}{a) $N(30) = 64 \cdot 1{,}25^{30} \approx 51\,699$ – die Aussage stimmt. \quad b) $N(20) = 120 \cdot 2^{20} = 125\,829\,120$ – die Aussage stimmt. \quad c) $x = 12$: $36 \cdot 1{,}04^{12} \approx 57{,}6$; Studie 2: $36 \cdot 1{,}5 = 54$ – Studie 1 sagt für $2033$ den höheren Verbrauch voraus. \quad d) $x = 20$: $120 \cdot 1{,}02^{20} \approx 178{,}3$; Studie 2: $120 \cdot 1{,}4 = 168$ – Studie 1 liegt höher. \quad e) $8 \cdot 0{,}8^{24} \approx 0{,}04\,\mathrm{mg}$ \quad f) $80 \cdot 0{,}75^{24} \approx 0{,}08\,\mathrm{mg}$ \quad g) $800 \cdot 0{,}88^4 \approx 480\,\mathrm{hPa}$ – die Behauptung stimmt. \quad h) $1\,000 \cdot 0{,}85^5 \approx 444$ Lux – die Behauptung stimmt. \quad i) $2027$: $654\,840$ ($2024$); $700\,679$ ($2025$); $749\,726$ ($2026$); $802\,207\,€$ ($2027$) \quad j) $2027$: $8\,487$; $8\,742$; $9\,004$; $9\,274$; $9\,552$ Einwohner ($2027$)}
\erg{45}{a) das Guthaben – gesucht ist die Zeit, bis der Wert doppelt so groß ist \quad b) $300$ \quad c) $3$ Jahre – der doppelte Startwert $100$ steht bei Jahr $3$}
\erg{46}{a) nach etwa $18$ Jahren – den doppelten Startwert $72$ auf der senkrechten Achse suchen, waagerecht zum Graphen, senkrecht zur Zeitachse \quad b) nach etwa $28$ Jahren – Zielwert $240$ auf der senkrechten Achse, waagerecht bis zum Graphen, senkrecht hinunter zur Zeitachse \quad c) nach etwa $3$ Stunden ($4{,}10\,\mathrm{mg}$ liegt am nächsten an der Hälfte $4{,}00\,\mathrm{mg}$) \quad d) nach etwa $2$ Tagen ($45\,\mathrm{mg}$ liegt näher an $40\,\mathrm{mg}$ als $33{,}75\,\mathrm{mg}$)}
